\section{Related Work}
\label{sec:related}

\para{Sensory representations in LLMs.}
A central question in NLP is whether text-only models learn anything resembling perceptual knowledge. \citet{marjieh2023large} showed that GPT-4 predicts human pairwise similarity judgments across six sensory modalities (pitch, loudness, color, consonants, taste, timbre), recovering known perceptual structures such as the color wheel and pitch spiral. \citet{lee2025exploring} evaluated 21 models on the \lancaster~\citep{lynott2020lancaster}, finding that top models achieve 85--90\% accuracy on sensory strength ratings. These studies establish that LLMs \emph{can} produce sensory judgments when prompted, but leave open whether sensory representations are activated spontaneously during processing.

\para{Explicit sensory prompting.}
\citet{wang2025words} introduced ``generative representations''---averaged hidden states across autoregressive generation---and showed that sensory cues (``imagine seeing/hearing'') steer LLM representations toward alignment with vision (DINOv2) and audio (BEATs) encoders. Their work demonstrates modality-specific activation under explicit prompting but does not test implicit activation from narrative objects. Our work directly addresses this gap: we compare implicit narrative context against explicit sensory description to determine whether activation is spontaneous or requires overt sensory cues.

\para{Probing for sensory and perceptual features.}
Linear probing is a well-established technique for testing what information is encoded in neural network representations~\citep{alain2017understanding, belinkov2022probing}. \citet{hicke2025zero} trained linear probes on BERT, RoBERTa, T5, GPT-2, and Qwen representations to predict sensory language scores, achieving $R^2$ up to 0.85 for concreteness prediction. \citet{ngo2024what} used contrastive probes to retrieve audio representations from text-only LM hidden states, finding above-chance generalization to unseen objects. \citet{merullo2025mouth} discovered internal phonetic representations in LLaMA, including a vowel chart resembling the IPA. These results motivate our use of linear probes to test for spontaneous sensory encoding.

\para{The embodiment gap.}
Not all evidence supports rich sensory encoding in text-only models. \citet{xu2025large} systematically compared ${\sim}4{,}400$ concepts between humans and LLMs, finding that similarity decreases from non-sensorimotor to sensory to motor domains. \citet{lee2024language} introduced the H-Test benchmark and found near-chance LLM performance on tasks requiring knowledge of language's physical manifestation. \citet{siedenburg2023language} reported only partial correlation between ChatGPT's ratings of musical instrument sounds and human judgments. These findings highlight genuine limitations, but focus on behavioral outputs rather than internal representations---a distinction our probing approach is designed to address.

\para{Sensory language in generation.}
\citet{hicke2025zero} conducted the largest comparison of sensory language use between 19 LLMs and humans, finding that all models differ significantly from human sensory language patterns. They observed that RLHF training may discourage sensory language use, and noted an inverse relationship: ``the better a model understands a sensory axis, the less likely it is to use it.'' This intriguing dissociation between internal representation and behavioral output further motivates our approach of probing hidden states directly rather than analyzing generated text.

\para{Sensorimotor norms and grounding.}
The \lancaster~\citep{lynott2020lancaster} provide human ratings for 39,707 words across 6 perceptual and 5 action dimensions, serving as a widely-used ground truth for sensory research. \citet{wu2026finetune} used representational similarity analysis to show that fine-tuning can bridge the embodiment gap, with sensorimotor improvements generalizing across languages. We use the Lancaster norms both for stimulus selection and for continuous sensory strength prediction, providing a direct link between human perceptual ratings and LLM internal representations.

\para{Positioning our work.}
Our study differs from prior work in three key ways. First, we test \emph{implicit} sensory activation---whether mentioning ``thunder'' in a narrative activates auditory representations without any sensory verb---rather than relying on explicit prompts. Second, we provide a comprehensive analysis of \emph{all six} sensory modalities (including interoception) within a single model, characterizing their geometric relationships. Third, we combine categorical probing with continuous prediction and subspace geometry analysis, offering a multi-faceted view of how sensory information is encoded in LLM representations.
